\section*{4 Saving Time}
In the previous section we took the length $4n$ branching programs of [11,12] as a starting point to analyze whether $m$ could be significantly reduced while still maintaining a linear amortized size. In this section we investigate the opposite question: namely, is $4n$ optimal? If we do not restrict the amortized size of our program, then every function has a branching program of length $n$ regardless of $m$. We will not only consider branching programs of linear amortized size, but the stricter model of permutation branching programs.
We focus on permutation branching programs in this section for two reasons. First, all our algorithms in the previous section can be converted to permutation branching programs. To see this, we resist our connection between register programs and $m$-catalytic branching programs, as proven in Observation7. Our observation in fact says something stronger, which is that $f_n$ can be computed by a family of permutation branching programs of width $2^{s(n)+1}$. This follows because we can choose to not fix $R_{\mathrm{out}}$ to be0, and instead have one source node corresponding to each initial configuration $\tau_1\ldots\tau_s,\tau^{\mathrm{out}}$; by Definition6 this source reaches the sink node labeled $\tau_1\ldots\tau_s,\tau^{\mathrm{out}}+f(x)$, which is the identity permutation when $f(x)=0$ and a fixed set of $2^{s(n)}$ transpositions otherwise. Thus as a corollary of [11,12] we get the following:
\paragraph{Theorem16.} Every function $f$ can be computed by a read-4 permutation branching program of width $2^{2^n-1}$ (or $2^{\binom n{\leq d}-1}$, where $d$ is the $\mathbb F_2$-degree of $p_f$).
Given the strength of these results, only a factor of4 away from the best conceivable amortized size, there is no reason a priori to believe that permutation branching programs are too weak to consider even better upper bounds. Indeed we will show that improvements are possible.
